\chapter{Three-Valued and Paraconsistent Logic}
\label{ch:nonclassical}

Classical logic assumes every proposition is exactly one of true or
false, and that a contradiction proves anything. Both assumptions
can be dropped, separately or together, and dropping them yields
logics that are still precise and still mechanically checkable —
just not classical. This chapter works through the smallest cases:
a third truth value standing for a gap in information, and a fourth
standing for a genuine glut.

\section{A Third Value: Kleene's Strong Logic}

\begin{experiment}{Excluded middle fails with a gap value}
Introduce a type with three values — true, false, and
\emph{unknown} — and define conjunction, disjunction, and negation
by Kleene's rule: an operation returns \texttt{unknown} whenever the
classical truth table would need to consult an unknown input to
decide the result, and returns a definite value whenever one input
already settles the answer regardless of the other (for instance,
\texttt{false and unknown} is \texttt{false}, since a false
conjunct settles the conjunction no matter what the other conjunct
is).
\begin{lean}
inductive K3 where
  | tt | ff | unknown
  deriving DecidableEq, Repr

def K3.knot : K3 → K3
  | tt => ff
  | ff => tt
  | unknown => unknown

def K3.kor : K3 → K3 → K3
  | tt, _ => tt
  | _, tt => tt
  | ff, ff => ff
  | _, _ => unknown

theorem excluded_middle_fails :
    ∃ p : K3, K3.kor p p.knot ≠ K3.tt := by
  exact ⟨K3.unknown, by
    decide ⟩
\end{lean}
Unlike the classical case, $p \lor \lnot p$ is not a validity here:
when $p$ is \texttt{unknown}, so is $\lnot p$, and the disjunction of
two unknowns is itself unknown, never forced up to \texttt{tt}.
\end{experiment}

\section{A Fourth Value: Gluts and the Logic of Paradox}

\texttt{K3}'s third value stood for a \emph{gap} — genuinely
unknown, neither true nor false. The dual notion is a \emph{glut}: a
proposition that is, in the sense a paraconsistent logic makes
precise, both true and false at once. Priest's Logic of Paradox
(LP) is built on exactly that idea. Rather than force gaps and
gluts into one four-valued lattice — a genuine bilattice, whose
connectives are easy to get subtly wrong without a way to check
them interactively — LP is presented here as a clean, self-contained
companion to \texttt{K3}, with its own three values and its own
notion of which values count as ``holding.''

\begin{experiment}{Excluded middle holds again, via a glut}
\begin{lean}
inductive LP where
  | t | f | b
  deriving DecidableEq, Repr

def LP.lnot : LP → LP
  | t => f
  | f => t
  | b => b

def LP.lor : LP → LP → LP
  | t, _ => t
  | _, t => t
  | b, b => b
  | b, f => b
  | f, b => b
  | f, f => f

def LP.designated : LP → Bool
  | t => true
  | b => true
  | f => false

theorem lp_excluded_middle :
    ∀ p : LP, (LP.lor p p.lnot).designated = true := by decide
\end{lean}
\texttt{b} (``both'') and \texttt{t} count as \emph{designated} —
holding, in the sense relevant to validity — while only \texttt{f}
does not. Where \texttt{K3}'s gap value made $p \lor \lnot p$ fail
(an unknown disjoined with an unknown stayed unknown), LP's glut
value does not: $b \lor \lnot b = b \lor b = b$, and $b$ is
designated. Excluded middle survives contact with a glut in exactly
the way it did not survive contact with a gap.
\end{experiment}

\section{Why Explosion Fails: A Countermodel}

\begin{countermodel}{Ex falso quodlibet is not valid in LP}
\begin{lean}
def LP.land : LP → LP → LP
  | f, _ => f
  | _, f => f
  | b, b => b
  | b, t => b
  | t, b => b
  | t, t => t

example : (LP.land LP.b (LP.lnot LP.b)).designated = true := by decide
example : LP.f.designated = false := by decide
\end{lean}
At $p = \texttt{b}$, $p \land \lnot p$ evaluates to \texttt{b} —
designated, so the premise counts as holding — while $q := \texttt{f}$
is not designated at all. Explosion, ``from a contradiction, anything
follows,'' would require every $q$ to be forced designated once
$p \land \lnot p$ is; this single witness is enough to show it is
not. A contradiction can hold in LP without licensing everything.
\end{countermodel}

LP is not a detour from this book's larger concerns; it is the
starting point for them. The constraint-regime logic developed at
length in the Priest–Flyxion program builds on exactly this
observation — that admitting some contradictions without admitting
everything is a coherent, checkable stance, not a confusion to be
regulated away — and \texttt{LP.designated} above is the simplest
possible instance of a \emph{constraint} on which values count as
holding, the same idea Chapter~\ref{ch:admissibility} develops in
full.

\section{Exercises}

\begin{exercise}
Define Kleene conjunction \texttt{K3.kand} and verify by
\texttt{decide} that \texttt{K3.knot} distributes over it the way De
Morgan's law requires.
\end{exercise}
